\section{Model interfaces and starter
architectures}\label{model-interfaces-and-starter-architectures}

The two launch interfaces, A for detection and B for transformation,
define compatible jobs and outputs. Miners may
change their internal architectures, grammatical representations, search
methods, and training data without changing those interfaces. Sharing a
foundation model between tasks is permitted. A detectors must be published
as complete immutable artifacts; an update enters a later qualification
round. B emission eligibility requires a public checkpoint nominated for
the monthly release cycle below. Miners may keep newer versions private
for training and paid hosting between releases. Training recipes and
data need not be public, subject to the rights required for each
released artifact. Publication terms must permit free artifact access,
local and commercial inference, redistribution, and reuse for training,
including private models, while retaining applicable license obligations.
A paid API alone does not qualify as publication. Hosting either task
may charge for compute and delivery.

\subsection{Author and origin
detection}\label{author-and-origin-detection}

An author task supplies a query and reference works for a gallery of
candidate authors. The response contains probabilities over the supplied
author IDs. The starter compares full sparse word/grammar profiles; a
learned text encoder with an author embedding head is the next
architectural comparison. Reference prototypes support new customer
authors without fitting a permanent classifier for each identity.
Unknown-author rejection requires a separate calibrated output and
withheld-author evaluation.

An origin task returns probabilities over documented production
histories: human-only, model-only, and mixed human/model work. The
production record includes revision order and assistance type. A
model-to-model rewrite remains model-only under this initial label
scheme. Pangram supplies a comparison, not the training label. The
proposed encoder uses a separate origin head; author similarity must not
be substituted for provenance.

ModernBERT-base~\cite{modernbert},
a 149M encoder, is a concrete starter candidate. Selection must compare
it with the existing metric on new authors, topics, and registers. The
subnet does not require that checkpoint or assume it will win.

\paragraph{Published artifact and execution policy.}\label{a-artifact-execution}
Each A artifact binds weights, features, tokenizer, preprocessing and
postprocessing, calibration, reference pooling, executable dependencies,
runtime, numerical precision, hardware requirements and output rules.
Pin prompts and random generators where used. Qualification must establish
a reproducible reference execution, including fixed tolerances and a
resolution rule for numerical boundary cases. Validate and cache the
qualified panel, its artifact hashes and thresholds before revealing the
B task or allowing candidate generation. Validators execute
that panel themselves for authoritative benchmark probabilities; a miner's
reported score is not the scoring oracle.

The panel uses distinct UIDs and distinct inference-content hashes,
excluding owner and signature metadata from the content identity. Identical
content receives one model-credit entry and one possible panel seat; the
earliest finalized accepted commitment wins, with canonical UID order for
ties. Small changes and equivalent implementations remain a qualification
problem. Select any reserve artifact before task issue. After issue, a
runner outage permits execution of the same artifact on a reserved runner,
not substitution of another detector or threshold. An unresolved reference
execution voids the common comparison under the closure rules; it is
neither B nonresponse nor an automatic pass.

Independent hidden author and production labels qualify each artifact.
Include targeted-trigger controls: a published model can faithfully execute
a rule designed to favor a colluding B. Publication verifies computation,
not reliable detection or independent ownership. The isolated runner
provides no network or undeclared state and excludes miner identities
and assignment metadata from model inputs; text itself
can still reveal task roles, so this does not prove blindness.

\subsection{Author transformation and detector
evasion}\label{author-transformation-and-detector-evasion}

Inputs are the source, independent target-author samples, requested
style changes, and a preservation brief. The brief identifies unwanted
writing patterns to remove and intentional voice features to retain.
Before generation, it fixes the target tonal intent using any requested
style and supplied reference samples. A dramatic authorial-style shift can change
tone while preserving the source's full meaning.
Author matching covers lexical, grammatical and rhetorical habits in
the requested register; generic polish alone does not satisfy it.
Every B task pursues author matching and evasion against Pangram and the
common subnet AI-origin panel defined in Section~\ref{purpose-and-scope}.
The miner returns exact revised text or an abstention. Every B benchmark
requires target-author references, style calibration and the complete
detector evidence. Opponent selection and scoring are fixed by the benchmark,
independently of the submitting miner. Author-directed revision of
one's own or licensed writing and authorized detector-robustness research
motivate this task. A detector pass does not replace the documented
production history or any disclosure required for the intended use.

The proposed starter is an 8B decoder, initially
Qwen3-8B~\cite{qwen}, with a style
adapter and a candidate-selection process. Named word and grammar
deviations can condition the first version. A learned profile prefix is
a later comparison, contingent on better results for unseen authors.
Deterministic Rust edit rules can propose small changes; the decoder can
propose broader restructuring.

Public A artifacts let B miners run local training queries without a
protocol query quota, at their own compute cost. They can use published
ensembles and retired A snapshots instead of depending on a peer's
availability for every search step. Direct gradients are available only
where the architecture permits them; differentiable feedback does not
guarantee useful discrete text edits. Training tickets serve authorized
rewrite requests from A to B under bounded limits. Benchmark submissions retain the frozen
candidate and deadline rules.
B can compute its own candidates' A scores before and after commitment;
these public-model scores are not hidden feedback. The official comparison
uses the one committed candidate and the frozen panel. Hidden labels,
other miners' candidates and preservation-review evidence remain restricted.

B submits the complete self-scored probability vectors and declared scalar
projections for every assigned A artifact, binding its hash, settings,
exact source inputs and final candidate. Validators independently run
those artifacts on the same inputs and compare the canonical outputs.
For emissions, they also reproduce the candidate with B's nominated
public artifact under the rules below. A miners need not serve inference
calls or help calculate these benchmark scores. Mandatory training-ticket
service and paid hosting may use newer private B versions; their outputs
cannot substitute for the published artifact's benchmark candidate.

Train first on reviewed source/reference/brief/revision tuples. Then
learn preferences among candidates that pass preservation review.
Include originals, failed rewrites, copied reference passages, and
fluent meaning-breaking edits as controls. Also include polished
revisions that erase the writer's voice and detector-passing revisions
that retain unwanted filler or formulaic phrasing. A miner's own style score may
guide its search, but cannot approve its output for benchmark rewards.

Transformation briefs may specify audience, clarity, accessibility,
organization, rhetorical force, and tone. A lower reading level or
shorter text is useful only when the brief calls for it. Suitable text
can remain unchanged. Validators use independent judges and reader
audits to assess preservation, readability, and requested voice;
reviewer ties and disagreements remain in the data.

Future comparisons of 4B, 8B, and 32B transformation models should use fidelity,
reader preference, and cost per accepted revision. All checkpoint
licenses and revisions must be pinned before a training release. The
named starters are proposals, not trained subnet products.

\subsection{Monthly B publication and local reference}\label{monthly-b-release}

The subnet maintains a downloadable B reference for customer-controlled
inference. Each UTC calendar month defines a nomination cycle. Before
the cycle begins, publish its nomination cutoff, qualification and
comparison windows, and future activation epoch as finalized block-height
boundaries. There is no obligation to replace a stronger incumbent merely
because a month has passed. At launch, a qualified published starter must
exist before advertising a local B service.

\paragraph{What miners publish.}
A B miner seeking emissions nominates one complete inference artifact
by the cutoff. It binds weights, or an adapter and an exactly pinned
downloadable base, tokenizer, prompts, preprocessing, author-profile
construction, edit/search/selection code, dependencies, numerical precision,
reference runtime and resource limits. Qualification requires the entire
package to run offline, without a secret key, remote API, telemetry or
undeclared state. Release rights cover every required component. A weight
file that depends on an unpublished selection service is incomplete.
Keep content-addressed copies and license notices in independent mirrors;
validators verify and cache the complete artifact before task disclosure.

Exact inference-content copies receive one B model-credit entry, using
the earliest finalized accepted commitment and canonical UID order for
ties, as for A. Publisher metadata does not create distinct content.
Repackaging or small weight changes do not prove an independent
contribution; evaluate their utility under the same budget. Republishing
the same qualifying checkpoint is allowed and supplies no novelty bonus.

\paragraph{Bind emissions to the released model.}
The nominated artifact stays fixed for its monthly scoring cycle.
The benchmark fixes a task seed independently of the submitter using
the protocol domain, cycle and shared comparison index. The reference
execution must produce exact final candidate text under the frozen
runtime, seed and compute budget, including any search and selection.
It may use the task's frozen public A panel locally within that budget;
no live Pangram call or other external service is part of generation.
B submits that candidate and its required detector evidence. Validators
independently reproduce the text before applying the common A/Pangram,
author-fit and semantic rules. A mismatch certified under the strict
weighted validity rule invalidates that submission and earns zero utility,
without substituting the replayed text. Runner disagreement or unavailable common evidence follows
the existing whole-comparison closure policy. Miners cannot substitute
text from stronger private weights, choose a different seed or make an
unrecorded final edit. Pangram reports concern the resulting fixed text;
their permitted repeat-selection policy remains explicit.

Publication and successful qualification are additional B eligibility
conditions, alongside the existing service and recovery-deficit gates.
Missing or incomplete publication earns no B skill for the affected
cycle. Newly nominated weights enter a later cycle. Certification and
native weight activation remain prospective under the actual reveal
lag; this rule promises no retroactive clawback. Paid demand and training
claims do not add emission credit.

\paragraph{Select the reference.}
Compare qualifying released candidates and the incumbent on the same
held-out source families, target-author briefs, frozen style evaluator,
A panel and Pangram policy. Use the joint B utility in
Equation~\ref{eq:transformation-utility}, with all preservation,
readability and target-tone gates, under equal declared compute limits.
Before task issue, freeze the qualification floors, aggregate ranking,
paired source-family uncertainty test and selection correction for the
number of challengers. Confirm a proposed replacement on independently
held-out families. Publish aggregate evidence, supported task limits,
resource requirements, artifact hashes and the future activation epoch
under the existing strict weighted certificate rule. Reserve a separate
fixed comparison budget, including generation replay and full semantic
review, before nomination opens. The operator funds incumbent reports
and the independent confirmation round, including when the incumbent's
publisher has left; ordinary candidate submissions retain the miner-funded
report policy. A retained artifact can participate as a reference without
granting emissions to an absent or ineligible publisher. Missing comparison
funding or evidence prevents promotion. Reference selection adds no reward pool.

Recommend the best independently confirmed released checkpoint under
that comparison. Retain the incumbent on ties, inconclusive comparisons
or no qualifying improvement. If the incumbent no longer qualifies,
suspend its recommendation until a replacement qualifies; retain its
archive and failure evidence. ``Best'' refers to evaluated public
submissions under the stated budget, not every private service or every
writer. Release scores cannot certify any particular customer's revision.

\paragraph{Private lead and public access.}
Miners may sell newer private versions between nomination cycles, with
separate service-version identifiers and performance evidence. Public
checkpoint scores do not attest which weights a hosted endpoint runs.
At the next cutoff, a miner may nominate its improved version to compete
for later emissions or retain it privately and continue competing with
its older qualified artifact. The protocol rewards demonstrated public
performance; it cannot compel disclosure of an unknown private model.

Monthly releases let customers run B without disclosing their text to a
miner. They also let Pangram, other detector providers and competing
miners download, train against and improve those snapshots. Only
unreleased improvements retain weight secrecy, and served outputs remain
observable to their recipients. The economic premise is that emissions
for released performance and revenue from newer hosted versions can
support continuing improvement. Measure that premise; neither a durable
private lead nor a hosted price premium is guaranteed.
